\chapter{Chest}

\label{ch:ch41_chest}

Chest imaging is used to investigate cough, breathlessness, chest pain, and abnormal blood tests. The most common investigations are the \textbf{chest X-ray} (CXR), which is quick and widely available, and \textbf{CT of the chest}, which provides much more detailed cross-sectional images. CT is particularly useful for characterising lung nodules, detecting blood clots in the lungs (pulmonary embolism), and staging cancers.

A chest CT may be performed \textbf{without contrast} (to look at the lung tissue and airways) or \textbf{with intravenous contrast} (to evaluate blood vessels, the mediastinum, and lymph nodes). Contrast dye is injected through a vein and may cause a warm sensation. The scan takes about 5--10 minutes and requires you to hold your breath several times.

Your radiologist will systematically examine every part of the chest image---not just the lungs but also the airways, blood vessels, heart, pleura (lining of the lungs), bones, and soft tissues.

---

\section{Approach}
Chest radiography, CT, CT angiography, MRI, ultrasound, and nuclear medicine answer different questions. Review the clinical indication, prior studies, respiratory status, contrast need, and radiation considerations before selecting a protocol.

\begin{table}[h]
\centering
\caption{Chest imaging modalities and indications}
\begin{tabular}{p{3cm} p{5cm} p{5cm}}
\hline
\textbf{Modality} & \textbf{Strengths} & \textbf{Common indications} \\
\hline
Chest X-ray & Fast, low dose, portable & Initial assessment of cough, pneumonia, pleural effusion, pneumothorax \\
CT chest (IV contrast) & Mediastinum, vessels, lymph nodes & Staging, PE (CTPA), mediastinal masses \\
CT chest (non-contrast) & Lung parenchyma, nodules, interstitial disease & Lung nodule follow-up, haemoptysis, interstitial lung disease \\
CT pulmonary angiography (CTPA) & Pulmonary arterial tree & Suspected pulmonary embolism \\
MRI chest & Soft tissue contrast, no radiation & Mediastinal masses, chest wall invasion, cardiac masses \\
PET-CT & Metabolic activity & Cancer staging, infection surveillance, solitary pulmonary nodule \\
Ultrasound & Bedside, no radiation & Pleural effusion (diagnosis and drainage guidance), chest wall masses \\
\hline
\end{tabular}
\end{table}

---

\section{Lung Parenchyma}

\subsection{Lung Nodules}
A pulmonary nodule is a rounded or irregular opacity $\leq 3$ cm in diameter, completely surrounded by lung parenchyma. Incidental lung nodules are extremely common (found in 25--50\% of CT scans).

\textbf{Key descriptors}:
\begin{itemize}
  \item \textbf{Size}: Largest diameter in mm (nodule $< 6$ mm is usually benign)
  \item \textbf{Morphology}: Solid, part-solid (subsolid), or ground-glass
  \item \textbf{Shape}: Smooth, lobulated, or spiculated (spiculated = higher malignancy risk)
  \item \textbf{Calcification}: Benign patterns: central, laminated, popcorn, or diffuse. Eccentric calcification does not exclude malignancy.
  \item \textbf{Location}: Upper lobe nodules carry higher malignancy risk
  \item \textbf{Growth}: doubling time assessment from prior imaging
\end{itemize}

\begin{pearlbox}[Fleischner Society guidelines]
For solid nodules $\geq 6$ mm in low-risk patients: 6--8 mm nodule requires CT at 6--12 months, then at 18--24 months if stable. Nodules $> 8$ mm require CT at 6--12 months, then at 18--24 months, or consider PET-CT / biopsy. Subsidiary nodules (part-solid or ground-glass) require longer follow-up (up to 5 years for ground-glass nodules).
\end{pearlbox}

\subsection{Lung-RADS}
Lung-RADS (Lung CT Screening Reporting and Data System) is used for lung cancer screening CTs. It categorises findings into:

\begin{table}[h]
\centering
\caption{Lung-RADS categories}
\begin{tabular}{p{2cm} p{4cm} p{3.5cm} p{3cm}}
\hline
\textbf{Category} & \textbf{Description} & \textbf{Action} & \textbf{Malignancy risk} \\
\hline
1 & Negative & Annual screening & $< 1\%$ \\
2 & Benign appearance & Annual screening & $< 1\%$ \\
3 & Probably benign & Short-interval CT (6 months) & $1$--$2\%$ \\
4A & Suspicious & CT in 3 months & $5$--$15\%$ \\
4B & Highly suspicious & Tissue sampling or PET-CT & $> 15\%$ \\
4X & Category 3 or 4 with additional suspicious features & Tissue sampling & High \\
\hline
\end{tabular}
\end{table}

\subsection{Consolidation and Ground-Glass Opacity}
\begin{itemize}
  \item \textbf{Consolidation}: Homogeneous increase in lung attenuation that obscures underlying vessels and airways. Air bronchograms may be present. Causes: pneumonia, pulmonary oedema, haemorrhage, alveolar proteinosis.
  \item \textbf{Ground-glass opacity (GGO)}: Increased lung attenuation that does NOT obscure underlying vessels. Causes: infection, interstitial lung disease, alveolar haemorrhage, adenocarcinoma (lepidic pattern).
  \item \textbf{Crazy-paving pattern}: GGO with superimposed interlobular septal thickening. Causes: alveolar proteinosis (classic), lipoid pneumonia, ARDS.
\end{itemize}

---

\section{Pleural Disease}
\begin{itemize}
  \item \textbf{Pleural effusion}: Blunting of the costophrenic angle on CXR; CT shows dependent fluid. Exudates ($> 30$ HU on CT) suggest malignancy, infection, or embolism; transudates ($< 15$ HU) suggest heart failure.
  \item \textbf{Pneumothorax}: Air in the pleural space; visible as a visceral pleural line on CT. Tension pneumothorax is a clinical diagnosis with mediastinal shift---do not wait for imaging.
  \item \textbf{Pleural thickening}: Diffuse or focal. Diffuse thickening with calcification is associated with asbestos exposure. Focal thickening requires exclusion of mesothelioma.
  \item \textbf{Empyema}: Loculated pleural fluid with enhancing pleura (split pleura sign) on CT. May require drainage.
\end{itemize}

---

\section{Mediastinal Masses}
The mediastinum is divided into compartments, and the location narrows the differential:

\begin{table}[h]
\centering
\caption{Mediastinal compartments and common masses}
\begin{tabular}{p{3cm} p{5cm} p{4.5cm}}
\hline
\textbf{Compartment} & \textbf{Contents} & \textbf{Common masses} \\
\hline
Anterior & Thymus, thyroid, great vessels, lymph nodes & Thymoma, teratoma/germ cell tumour, thyroid mass, lymphoma \\
Middle & Heart, pericardium, great vessels, lymph nodes & Lymphadenopathy, bronchogenic cyst, pericardial cyst \\
Posterior & Oesophagus, descending aorta, spine, nerves & Neurogenic tumour (schwannoma, neurofibroma), oesophageal duplication cyst \\
\hline
\end{tabular}
\end{table}

\begin{pearlbox}[The 4 T's of anterior mediastinal masses]
The mnemonic ``4 T's'' helps remember the differential for an anterior mediastinal mass: \textbf{T}hymoma, \textbf{T}eratoma (germ cell tumour), \textbf{T}hyroid (goitre/substernal), and ``\textbf{T}errible'' lymphoma.
\end{pearlbox}

---

\section{Chest Wall}
Assess the ribs, sternum, thoracic vertebrae, muscles, and subcutaneous tissues. Rib fractures are common in trauma; pathological fractures suggest metastatic disease. Chest wall tumours include desmoid tumours, sarcomas, and metastatic deposits.

---

\section{Structured Report: Solitary Pulmonary Nodule}

\begin{testbox}[Indications:]
62-year-old male smoker. Incidental 12 mm solid nodule in the right upper lobe on staging CT for colorectal cancer. No prior chest imaging available.
\end{testbox}
\begin{refbox}[Findings:]
\textbf{Lung windows}: A 12 mm solid nodule is identified in the right upper lobe, segment 2. The nodule has mildly spiculated margins and no internal calcification. No additional pulmonary nodules. No ground-glass opacities. Lung parenchyma is otherwise clear.\\
\textbf{Mediastinal windows}: No mediastinal or hilar lymphadenopathy. No pleural effusion. No chest wall abnormality.
\end{refbox}
\begin{pearlbox}[Impression:]
1. 12 mm spiculated solid nodule, right upper lobe---suspicious for primary lung malignancy given spiculated morphology and smoking history.
2. Recommend PET-CT for further characterisation. If PET-CT is inconclusive or highly suspicious, recommend CT-guided biopsy or surgical resection.
3. Ensure prior chest imaging is obtained for comparison if available.
\end{pearlbox}

---

\section{Systematic Review}
Assess tubes and lines, lungs, airways, pleura, mediastinum, heart, vessels, diaphragm, bones, and soft tissues. Describe distribution, acuity, and complications rather than relying on a single sign.

\section{Communication}
Tension pneumothorax, major airway or line malposition, pulmonary embolism, aortic injury, severe edema, and large pleural collections may require urgent communication.

\clearpage
